% 제3장 연구 방법
% 1차 심사 반영: 구 2.12 연구 설계 + 구 3.1 시뮬레이션 설계 를 통합

\chapter{\vrev{연구 방법}}
\label{ch:method}

\vrev{%
본 연구의 연구 설계와 시뮬레이션 실험 방법을 기술하기 위하여 먼저
관측성과 불확실성의 조합으로 정의되는 4단계 조건 체계(C0--C3)와
이에 대응하는 \jrev{10개 가설}의 검증 설계를 제시한다.
이어서 $n = 10{,}000$ 몬테카를로 시뮬레이션의 실행 절차와
분포 매개변수, 시드 배정 규칙, 통계량 산출 방법을 제시한다.
}

\chg{본 연구는 모델 정립에서 산업 범용성 검증에 이르기까지 일곱 단계의 절차를
거치며, \figref{fig:comprehensive_framework}\refpo{fig:comprehensive_framework}{은}{는}
이 전체 절차를 한눈에 보여준다. 본 연구는 먼저 관측성과 불확실성을 시그모이드
바운딩으로 결합하는 동적 복구목표시간(DRTO) 수식 모델 2.7(DRTO 모델 구조)를 정립하는 데에서
출발한다. 다음으로 정적 기준에서 관측성과 불확실성을 단계적으로 도입하는 네 개
운영 조건(C0--C3)을 설계하여, 각 요인의 기여와 결합 효과를 분리하여 관찰할 수
있는 실험 틀을 마련한다. 본격적인 가설 검증에 앞서 기반 검증(L0)에 해당하는
Golden Compare를 수행하여, 수식의 코드 구현이 단일 기준 정의(SSOT)와 정확히
일치함을 결정론적으로 확인한다. 이 토대 위에서 검증은 네 층위로 이어진다.
$n=10{,}000$ 몬테카를로 시뮬레이션으로 관측성 효과와 경계 보장 등 모델의 기본
작동을 확률적으로 확인하고(L1), 2차 다항 회귀와 분할 회귀로 변수 간 설명력과
이중채널 비선형 상호작용을 규명하며(L2와 L4), 다중 시드 분석으로 결과가 특정
난수에 의존하지 않고 재현됨을 보인다(L3). 이어 BCM 및 재해복구 분야
전문가 평가로 모델의 이론적 타당성과 실무 적용 가능성을 삼각검증하고(L5),
마지막으로 여섯 개 산업 시나리오에 적용하여 모델이 특정 환경에 국한되지 않는
범용성을 갖는지 확인한다. 이처럼 각 단계는 앞 단계의 결과를 전제로 다음 단계로
이어지며, 하위 검증의 통과가 상위 검증의 신뢰성을 뒷받침하는 계층적 구조를
이룬다.}

% 제3장 도입 — 연구 전체 수행 절차 개관 그림 (1장에서 이동, 결과수치 제거됨)
\begin{figure}[htbp]
  \centering
  \resizebox{\textwidth}{!}{%
  \begin{tikzpicture}[
    >=Stealth,
    stepnum/.style={circle, draw=black, fill=white, font=\scriptsize\bfseries,
                    minimum size=5.5mm, inner sep=0pt, line width=0.7pt},
    mainbox/.style={draw=black, rounded corners=5pt, line width=0.7pt,
                    align=center, font=\scriptsize, minimum height=8mm,
                    inner sep=3pt},
    subbox/.style={draw=black, rounded corners=2pt, line width=0.4pt,
                   align=center, font=\tiny, minimum height=7mm,
                   minimum width=22mm, inner sep=2pt},
    condbox/.style={draw=black, rounded corners=2pt, line width=0.5pt,
                    align=center, font=\tiny, minimum width=16mm,
                    minimum height=8mm, inner sep=2pt},
    mainarr/.style={->, line width=1.0pt, black},
    subarr/.style={->, line width=0.5pt, black},
  ]

  %% 행 좌표 (상→하)
  \def\rowA{14.0}   % ① 수식 모델
  \def\rowB{11.5}   % ② 운영 조건       (①-②: 2.5)
  \def\rowC{7.5}    % ③ 시뮬레이션      (②-③: 4.0, ②-L0 화살표 공간 확보)
  \def\rowD{3.8}    % ④ 회귀·교차분석   (③-④: 3.7, H3 하위 2단 + 박스 간섭 회피)
  \def\rowE{1.6}    % ⑤ 강건성 검증     (④-⑤: 2.2)
  \def\rowF{-0.6}   % ⑥ 전문가 평가     (⑤-⑥: 2.2)
  \def\rowG{-2.8}   % ⑦ 산업 범용성     (⑥-⑦: 2.2)

  %% ====================================================================
  %% ① 수식 모델
  %% ====================================================================
  \node[stepnum] (n1) at (-1.3, \rowA) {1};
  \node[mainbox, fill=blue!8, minimum width=115mm] (m1) at (5.0, \rowA)
    {\textbf{수식 모델 정립}\;
    {\tiny $\DRTO = R_0 + E_{O,\mathrm{bnd}} + E_{U,\mathrm{bnd}}$,\;
    $S(z) = 2/(1\!+\!e^{-z})\!-\!1$,\; $\kappa = 1.2$}};

  \node[subbox, fill=blue!5] (m1a) at (1.3, \rowA-1.1)
    {관측성 채널\\$E_{O,\mathrm{bnd}} = -R_0 \!\cdot\! S(z_O)$};
  \node[subbox, fill=blue!5] (m1b) at (5.0, \rowA-1.1)
    {불확실성 채널\\$E_{U,\mathrm{bnd}} = (R_{\max}\!-\!R_0)\!\cdot\! S(z_U)$};
  \node[subbox, fill=blue!5] (m1c) at (8.7, \rowA-1.1)
    {경계 보장\\$0 < \DRTO < R_{\max}$};

  \draw[subarr] (m1.south -| m1a) -- (m1a.north);
  \draw[subarr] (m1.south -| m1b) -- (m1b.north);
  \draw[subarr] (m1.south -| m1c) -- (m1c.north);

  %% ====================================================================
  %% ② 운영 조건
  %% ====================================================================
  \node[stepnum] (n2) at (-1.3, \rowB+0.5) {2};
  \node[mainbox, fill=green!8, minimum width=115mm] (m2) at (5.0, \rowB+0.5)
    {\textbf{4단계 운영 조건 설계 (C0--C3)}};

  \node[condbox, fill=gray!15] (c0) at (0.8, \rowB-0.7)
    {\textbf{C0}\\정적 기준};
  \node[condbox, fill=blue!10] (c1) at (3.6, \rowB-0.7)
    {\textbf{C1}\\관측성};
  \node[condbox, fill=green!10] (c2) at (6.4, \rowB-0.7)
    {\textbf{C2}\\Gaussian};
  \node[condbox, fill=red!10] (c3) at (9.2, \rowB-0.7)
    {\textbf{C3}\\Pareto};

  \draw[subarr] (m2.south -| c0) -- (c0.north);
  \draw[subarr] (m2.south -| c1) -- (c1.north);
  \draw[subarr] (m2.south -| c2) -- (c2.north);
  \draw[subarr] (m2.south -| c3) -- (c3.north);

  \draw[subarr] (c0) -- (c1) node[midway, above, font=\tiny] {$+O$};
  \draw[subarr] (c1) -- (c2) node[midway, above, font=\tiny] {$+U^{(G)}$};
  \draw[subarr] (c1.south) -- ++(0,-0.7) -| (c3.south)
    node[pos=0.25, below, font=\tiny] {$+U^{(P)}$};

  %% ====================================================================
  %% L0: Golden Compare (② → ③ 사이)
  %% ====================================================================
  \node[stepnum, fill=black!8] (n0) at (-1.3, {(\rowB-1.5+\rowC)/2})
    {\tiny L0};
  \node[mainbox, fill=black!6, minimum width=115mm] (m0)
    at (5.0, {(\rowB-1.5+\rowC)/2})
    {\textbf{L0: 기반 검증 (Golden Compare)}\;
    {\tiny 47개 포인트 --- 순서 제약(H3), 물리적 경계(H2),
    $z_O$ 동일성, P99 비율 검증}};

  %% ====================================================================
  %% ③ 시뮬레이션 [L1: H1-H3] — H1~H3 각각 독립 박스, H3 하위 a/b/c
  %% ====================================================================
  \node[stepnum] (n3) at (-1.3, \rowC) {3};
  \node[mainbox, fill=yellow!10, minimum width=115mm] (m3) at (5.0, \rowC)
    {\textbf{몬테카를로 시뮬레이션}\; ($n\!=\!10{,}000 \times 4$)\;
    {\tiny [L1: H1--H3]}};

  \node[subbox, fill=yellow!6, minimum width=22mm] (m3h1) at (1.5, \rowC-1.1)
    {H1 관측성 효과};
  \node[subbox, fill=yellow!6, minimum width=22mm] (m3h2) at (5.0, \rowC-1.1)
    {H2 경계 보장};
  \node[subbox, fill=yellow!6, minimum width=28mm] (m3h4) at (8.7, \rowC-1.1)
    {H3 조건 순서};

  \draw[subarr] (m3.south -| m3h1) -- (m3h1.north);
  \draw[subarr] (m3.south -| m3h2) -- (m3h2.north);
  \draw[subarr] (m3.south -| m3h4) -- (m3h4.north);

  %% H3 하위 박스 (a, b, c)
  \node[subbox, fill=yellow!4, minimum width=12mm, minimum height=6mm]
    (h4a) at (7.4, \rowC-2.0)
    {a: $C1\!\leq\!C0$};
  \node[subbox, fill=yellow!4, minimum width=12mm, minimum height=6mm]
    (h4b) at (8.7, \rowC-2.0)
    {b: $C1\!\leq\!C2$};
  \node[subbox, fill=yellow!4, minimum width=12mm, minimum height=6mm]
    (h4c) at (10.0, \rowC-2.0)
    {c: $C2\!\leq\!C3$};

  \draw[subarr] (m3h4.south -| h4a) -- (h4a.north);
  \draw[subarr] (m3h4.south) -- (h4b.north);
  \draw[subarr] (m3h4.south -| h4c) -- (h4c.north);

  %% ====================================================================
  %% ④ 회귀·교차분석 [L2: H4-H5, L4: H8]
  %% ====================================================================
  \node[stepnum] (n4) at (-1.3, \rowD+0.5) {4};
  \node[mainbox, fill=orange!10, minimum width=115mm] (m4) at (5.0, \rowD+0.5)
    {\textbf{회귀분석 및 교차분석}\;
    {\tiny [L2: H4--H5, L4: H8]}};

  \node[subbox, fill=orange!6] (m4a) at (1.3, \rowD-0.6)
    {H4 회귀 설명력};
  \node[subbox, fill=orange!6] (m4b) at (5.0, \rowD-0.6)
    {H5 교호작용};
  \node[subbox, fill=orange!6] (m4c) at (8.7, \rowD-0.6)
    {H8 이중채널 감쇠};

  \draw[subarr] (m4.south -| m4a) -- (m4a.north);
  \draw[subarr] (m4.south -| m4b) -- (m4b.north);
  \draw[subarr] (m4.south -| m4c) -- (m4c.north);

  %% ====================================================================
  %% ⑤ 강건성 검증 [L3: H6-H7]
  %% ====================================================================
  \node[stepnum] (n5) at (-1.3, \rowE+0.5) {5};
  \node[mainbox, fill=cyan!8, minimum width=115mm] (m5) at (5.0, \rowE+0.5)
    {\textbf{강건성 검증}\;
    {\tiny [L3: H6--H7]}};

  \node[subbox, fill=cyan!5] (m5a) at (1.3, \rowE-0.6)
    {H6 시드 독립성};
  \node[subbox, fill=cyan!5] (m5b) at (5.0, \rowE-0.6)
    {H7 다중 시드 강건성};
  \node[subbox, fill=cyan!5] (m5c) at (8.7, \rowE-0.6)
    {Bonferroni 보정};

  \draw[subarr] (m5.south -| m5a) -- (m5a.north);
  \draw[subarr] (m5.south -| m5b) -- (m5b.north);
  \draw[subarr] (m5.south -| m5c) -- (m5c.north);

  %% ====================================================================
  %% ⑥ 전문가 평가 + 가설 종합 [L5: H9-H10]
  %% ====================================================================
  \node[stepnum] (n6) at (-1.3, \rowF+0.5) {6};
  \node[mainbox, fill=violet!8, minimum width=115mm] (m6) at (5.0, \rowF+0.5)
    {\textbf{다층적 타당성 검증}\;
    {\tiny [L5: H9--H10]}};

  \node[subbox, fill=violet!5] (m6a) at (1.3, \rowF-0.6)
    {H9 전문가 정량 평가};
  \node[subbox, fill=violet!5] (m6b) at (5.0, \rowF-0.6)
    {H10 전문가 정성 평가};
  \node[subbox, fill=violet!5] (m6c) at (8.7, \rowF-0.6)
    {가설 종합 판정};

  \draw[subarr] (m6.south -| m6a) -- (m6a.north);
  \draw[subarr] (m6.south -| m6b) -- (m6b.north);
  \draw[subarr] (m6.south -| m6c) -- (m6c.north);

  %% ====================================================================
  %% ⑦ 산업 범용성 검증
  %% ====================================================================
  \node[stepnum] (n7) at (-1.3, \rowG+0.5) {7};
  \node[mainbox, fill=teal!8, minimum width=115mm] (m7) at (5.0, \rowG+0.5)
    {\textbf{6개 산업 시나리오 범용성 검증}};

  \node[subbox, fill=teal!4, minimum width=16mm] (m7a) at (0.0, \rowG-0.7)
    {제조업};
  \node[subbox, fill=teal!4, minimum width=16mm] (m7b) at (2.0, \rowG-0.7)
    {금융};
  \node[subbox, fill=teal!4, minimum width=16mm] (m7c) at (4.0, \rowG-0.7)
    {의료};
  \node[subbox, fill=teal!4, minimum width=16mm] (m7d) at (6.0, \rowG-0.7)
    {IT서비스};
  \node[subbox, fill=teal!4, minimum width=16mm] (m7e) at (8.0, \rowG-0.7)
    {물류};
  \node[subbox, fill=teal!4, minimum width=16mm] (m7f) at (10.0, \rowG-0.7)
    {소매};

  \draw[subarr] (m7.south -| m7a) -- (m7a.north);
  \draw[subarr] (m7.south -| m7b) -- (m7b.north);
  \draw[subarr] (m7.south -| m7c) -- (m7c.north);
  \draw[subarr] (m7.south -| m7d) -- (m7d.north);
  \draw[subarr] (m7.south -| m7e) -- (m7e.north);
  \draw[subarr] (m7.south -| m7f) -- (m7f.north);

  %% ====================================================================
  %% 왼쪽: 단계 진행 화살표 (흐름 명시)
  %% ====================================================================
  \foreach \a/\b in {n1/n2, n2/n0, n0/n3, n3/n4, n4/n5, n5/n6, n6/n7} {
    \draw[->, black, line width=0.7pt] (\a.south) -- (\b.north);
  }

  %% ====================================================================
  %% 오른쪽: 중괄호 (왼쪽 열림 = 본문 쪽) + 가설 매핑
  %% ====================================================================
  \draw[decorate, decoration={brace, amplitude=4pt},
        line width=0.5pt, black]
    (11.0, \rowA+0.4) -- (11.0, \rowB-1.0)
    node[midway, right=5pt, font=\tiny, text=black, align=left]
    {이론적\\기반};

  \draw[decorate, decoration={brace, amplitude=4pt},
        line width=0.5pt, black]
    (11.0, \rowC+0.35) -- (11.0, \rowC-2.35)
    node[midway, right=5pt, font=\tiny, text=black, align=left]
    {L1\\시뮬레이션};

  \draw[decorate, decoration={brace, amplitude=4pt},
        line width=0.5pt, black]
    (11.0, \rowD+0.9) -- (11.0, \rowD-1.0)
    node[midway, right=5pt, font=\tiny, text=black, align=left]
    {L2--L4\\회귀$\cdot$이중채널};

  \draw[decorate, decoration={brace, amplitude=4pt},
        line width=0.5pt, black]
    (11.0, \rowE+0.9) -- (11.0, \rowE-1.0)
    node[midway, right=5pt, font=\tiny, text=black, align=left]
    {L3\\강건성};

  \draw[decorate, decoration={brace, amplitude=4pt},
        line width=0.5pt, black]
    (11.0, \rowF+0.9) -- (11.0, \rowF-1.0)
    node[midway, right=5pt, font=\tiny, text=black, align=left]
    {L5\\전문가};

  \draw[decorate, decoration={brace, amplitude=4pt},
        line width=0.5pt, black]
    (11.0, \rowG+0.9) -- (11.0, \rowG-1.0)
    node[midway, right=5pt, font=\tiny, text=black, align=left]
    {적용};

  %% 바운딩 박스 확장 (우측 라벨 잘림 방지, 하단 확장)
  \path (13.5, \rowG-1.5);

  \end{tikzpicture}%
  }% end resizebox
  \caption{\jrev{본 연구의 전체 수행 절차 개관}}
  \label{fig:comprehensive_framework}
\end{figure}
    % 연구 전체 수행 절차 개관 그림 (1장에서 이동)



%% ===================================================================
\section{\vrev{연구 설계}}
\label{sec:research-design}
%% ===================================================================

%% -------------------------------------------------------------------
\subsection{\vrev{6단계 다층 검증 프레임워크}}
\label{subsec:validation-framework}
%% -------------------------------------------------------------------

\jrev{10개 가설}을 체계적으로 검증하기 위해 6단계 다층 검증(six-stage multi-layer
validation) 프레임워크를 설계한다. 각 단계는 서로 다른 검증 관점을 제공하며,
하위 단계의 결과가 상위 단계의 전제 조건으로 작용하는
계층적 의존 구조를 갖는다~\citep{oberkampf2010verification}.

\chg{프레임워크는 결정론적 기반 검증에서 출발하여 통계적 검증과 전문가 평가로 단계적으로 심화한다. 먼저 기반 검증(L0, Foundation Verification)은 Golden Compare로서, 7개 산업 시나리오의 대표 매개변수를 $\DRTO$ 수식에 직접 대입하여 해석적 기대값과 코드 출력을 비교하는 결정론적 검증이다. 순서 제약, 물리적 경계, $z_O$ 동일성, P99 비율 일관성의 네 범주에 걸친 총 47개 검증 포인트로 구성되며, 경계 보장(H2)과 순서 보존(\jrev{H3})을 결정론적으로 사전 검증한다. 47개 포인트 전체에서 PASS를 달성해야만 이후 L1에서 L5까지의 가설 검증이 의존하는 수식 구현이 단일 기준 정의(SSOT)와 정합함이 보장된다(결과는 \subsecref{subsec:golden_compare} 참조, 7개 시나리오의 전체 계산과 47개 점별 결과는 공개 재현 패키지에 수록).}

\chg{이 기반 위에서 시뮬레이션 검증(L1)은 $n = 10{,}000$ 몬테카를로 시뮬레이션으로 관측성 효과(H1), 경계 보장(H2), 조건 순서(\jrev{H3})를 확인하여 모델의 기본 행동을 검증한다. 이때 H2와 \jrev{H3}는 L0에서 결정론적으로 사전 확인된 뒤 본 단계의 확률적 표본으로 재검증되며, 곡률 매개변수 $\kappa = 1.2$의 최적성은 \citet{lee2026drto}에서 확정된 전제로서 본 연구의 가설 검증 대상이 아니다. 이어 회귀 검증(L2)은 2차 다항 회귀 모델의 설명력(\jrev{H4})과 교호작용 유의성(\jrev{H5})을 통해 응답 곡면의 구조적 특성을 규명하고, 강건성 검증(L3)은 시드 독립성(\jrev{H6})과 다중 시드 강건성(\jrev{H7})을 통해 결과의 재현성과 대표성을 확인하는데, 표본 $n = 10{,}000$과 시드 $10{,}000$개를 곱한 총 $10^8$회 계산으로 통계적 안정성을 보장한다. 다음으로 이중채널 검증(L4)은 P85.8 기준 분할 회귀로 이중채널 비선형 감쇠 효과의 존재와 크기를 검증하며, L1에서 L3까지의 결과를 종합하는 핵심 발견으로서 실무적 정책 시사점의 근거를 제공한다. 마지막으로 전문가 평가(L5)는 BCM 및 재해복구 전문가 5명의 정량적 설문(\jrev{H9})과 정성적 인터뷰(\jrev{H10})를 통해 모델의 이론적 타당성과 실무 적용 가능성을 삼각검증한다.}

\chg{특히 최상위 단계인 L5의 전문가 평가는 본 연구가 삼각검증(triangulation)을 채택한 지점이다. 시뮬레이션과 회귀(L1에서 L4까지)가 제공하는 정량적 증거만으로는 모델이 실제 업무연속성 현장에서 타당한지를 온전히 확증하기 어렵다. 단일 방법에 내재한 한계를 보완하기 위해, 정량적 분석 결과를 전문가의 정성적 판단과 교차 대조하여 결론이 서로 수렴하는지를 확인하는 삼각검증~\citep{creswell2018designing}을 결과 단계가 아니라 설계 단계에서부터 검증 체계의 최상위 층위로 포함한다. 이로써 정량적으로 도출된 모델의 효과가 현장 전문가의 실무적 판단과 일치할 때 비로소 타당성이 확증되는 구조를 갖춘다.}

\p{[}그림~\vref{fig:validation-layers}\p{]}는 6단계 검증 프레임워크의
구조와 계층 간 의존 관계를 도식화한 것이다.

\begin{figure}[htbp]
\centering
\resizebox{0.95\textwidth}{!}{%
\begin{tikzpicture}[
  pyr/.style={
    draw, rounded corners=5pt, minimum height=1.3cm,
    align=center, font=\small, line width=0.6pt,
    drop shadow={shadow xshift=0.8mm, shadow yshift=-0.8mm, opacity=0.15}
  },
  arr/.style={-{Stealth[length=2.5mm]}, thick, black},
  lnum/.style={circle, draw, fill=white, font=\small\bfseries, inner sep=2pt,
    minimum size=7mm, line width=0.6pt},
  rlbl/.style={font=\scriptsize, text=black, align=left, anchor=west},
]

% 피라미드 계층 (아래→위, 넓은→좁은)
\node[pyr, fill=black!8, minimum width=15.5cm] (L0) at (0, -1.8) {
  \textbf{L0: 기반 검증 (Golden Compare)} \quad
  47개 포인트 \;\textbar\;
  순서 제약 \;\textbar\;
  물리적 경계 \;\textbar\;
  $z_O$ 동일성 \;\textbar\;
  P99 비율
};
\node[pyr, fill=green!15, minimum width=14cm] (L1) at (0, 0) {
  \textbf{L1: 시뮬레이션 검증} \quad
  H1 (관측성 효과) \;\textbar\;
  H2 (경계 보장) \;\textbar\;
  H3 (순서 보존)
};
\node[pyr, fill=violet!10, minimum width=12cm] (L2) at (0, 1.8) {
  \textbf{L2: 회귀 검증} \quad
  H4 ($R^2_{(2\mathrm{q})} \geq 0.95$) \;\textbar\;
  H5 (교호작용)
};
\node[pyr, fill=blue!10, minimum width=10cm] (L3) at (0, 3.6) {
  \textbf{L3: 강건성 검증} \quad
  H6 (시드 독립성, $|r| < 0.02$) \;\textbar\;
  H7 (다중 시드, $10^8$회)
};
\node[pyr, fill=orange!15, minimum width=8cm] (L4) at (0, 5.4) {
  \textbf{L4: 이중채널}\\
  H8 (감쇠비 $< 1$)
};
\node[pyr, fill=yellow!10, minimum width=6cm] (L5) at (0, 7.2) {
  \textbf{L5: 전문가 평가}\\
  H9 ($\bar{X} \geq 4.0$) \;\textbar\;
  H10 (동의율 $\geq 80\%$)
};

% 계층 번호 (왼쪽 원형 배지)
\node[lnum, text=black] at ([xshift=-9.1cm]L0.center) {L0};
\node[lnum, text=black] at ([xshift=-8.4cm]L1.center) {L1};
\node[lnum, text=black] at ([xshift=-6.8cm]L2.center) {L2};
\node[lnum, text=black] at ([xshift=-6.4cm]L3.center) {L3};
\node[lnum, text=black] at ([xshift=-4.8cm]L4.center) {L4};
\node[lnum, text=black] at ([xshift=-3.8cm]L5.center) {L5};

% 우측 핵심 산출물
\node[rlbl] at ([xshift=8.5cm]L0.center) {SSOT 구현 정합성};
\node[rlbl] at ([xshift=7.8cm]L1.center) {$n=10{,}000$ 시뮬레이션};
\node[rlbl] at ([xshift=6.8cm]L2.center) {응답 곡면 규명};
\node[rlbl] at ([xshift=5.8cm]L3.center) {재현성 $10^8$회};
\node[rlbl] at ([xshift=4.8cm]L4.center) {이중채널 효과 검증};
\node[rlbl] at ([xshift=3.8cm]L5.center) {전문가 삼각검증};

% 계층 간 화살표
\draw[arr] (L0.north) -- (L1.south);
\draw[arr] (L1.north) -- (L2.south);
\draw[arr] (L2.north) -- (L3.south);
\draw[arr] (L3.north) -- (L4.south);
\draw[arr] (L4.north) -- (L5.south);

% 좌측 방향 표시
\node[font=\footnotesize, text=black, rotate=90] at (-9.8, 2.7) {검증 심화 방향};
\draw[-{Stealth[length=2mm]}, thick, black] (-9.4, -2.3) -- (-9.4, 7.7);

\end{tikzpicture}%
}% end resizebox
\caption{6단계 다층 검증 프레임워크}
\label{fig:validation-layers}
\end{figure}

\p{[}표~\ref{tab:six-stage}\p{]}은 6단계 검증 프레임워크의 각 단계와 대응 가설,
검증 목적을 요약한 것이다. 여기서 ``L''은 단계(Level)의 약자로, L0의 결정론적
구현 정합성 검증에서 출발하여 L5의 전문가 삼각검증에 이르는 누적적 위계를 구성하며,
각 단계는 이전 단계의 결과를 전제로 한다. 이 순서는 본 연구의 실행 흐름과 일치한다.

\begin{table}[H]
\centering
\caption{6단계 다층 검증 프레임워크}
\label{tab:six-stage}
\small
\begin{tabular}{c l c l}
\toprule
\textbf{단계} & \textbf{명칭} & \textbf{가설} & \textbf{목적} \\
\midrule
L0 & 기반 검증       & (H2, \jrev{H3} 사전) & Golden Compare 47개 포인트: SSOT 구현 정합성 \\
L1 & 시뮬레이션 검증 & H1--\jrev{H3}   & 모델 기본 행동 확인 \\
L2 & 회귀 검증       & \jrev{H4}--\jrev{H5}   & 설명력과 교호작용 검증 \\
L3 & 강건성 검증     & \jrev{H6}--\jrev{H7}   & 재현성과 안정성 보장 \\
L4 & 이중채널 검증   & \jrev{H8}       & 핵심 발견 규명 \\
L5 & 전문가 평가     & \jrev{H9}--\jrev{H10} & 실무 타당성 삼각검증 \\
\bottomrule
\end{tabular}
\end{table}

이 위계에서 L0의 기반 검증과 L1 이후의 통계적 검증은 성격이 근본적으로 다르다.
L0의 결정론적 검증(deterministic verification)은 고정된 입력 매개변수에서 모델 수식의
해석적 출력과 코드 실행 결과를 직접 비교하여 정확 일치 여부를 PASS/FAIL로 판정하는
방식으로, 확률적 표본 추출이나 통계적 추론을 사용하지 않으며 동일 입력에서 항상 동일
출력이 보장된다. 반면 L1--L4의 통계적 검증은 확률 분포에서 $n = 10{,}000$ 표본을
추출하여 평균, 분산, 백분위, 회귀 계수 등 모델 거동의 통계적 성질을 분석한다.
두 방식의 차이를 \p{[}표~\ref{tab:verification-types}\p{]}에 정리하였다.

\begin{table}[H]
\centering
\caption{결정론적 검증과 통계적 검증의 비교}
\label{tab:verification-types}
\small
\begin{tabular}{l p{5.5cm} p{5.5cm}}
\toprule
\textbf{항목} & \textbf{결정론적 검증 (L0)} & \textbf{통계적 검증 (L1--L4)} \\
\midrule
입력 & 고정 매개변수 (시나리오) & 확률 분포 ($n = 10{,}000$ 표본) \\
출력 비교 & 해석적 기대값 = 코드 출력 (정확 일치) & 분포·평균·검정 통계량 \\
결론 형식 & PASS / FAIL (이진) & 가설 채택/기각, 효과크기 \\
변동성 & 없음 (재실행 시 동일) & 시드·표본에 따른 변동 \\
검증 대상 & 수식 구현 정합성 (코드와 이론) & 모델 거동의 통계적 성질 \\
대표 사례 & Golden Compare 47개 포인트 & Cohen's $d$, $R^2$, CVaR \\
\bottomrule
\end{tabular}
\end{table}

두 방식은 위계적 사전 조건 관계를 갖는다. 결정론적 검증(L0)이 PASS되어야 수식이
코드로 정확히 구현됨이 보장되며, 비로소 후속 통계적 검증(L1--L4)의 결과가 의미를
갖는다. 만약 L0에서 단 하나의 검증 포인트라도 FAIL이면 코드 구현 자체에 오류가
있다는 의미이므로 이후의 모든 분석 결과를 신뢰할 수 없게 된다~\citep{oberkampf2010verification}.
이러한 ``결정론적 사전 검증에서 통계적 본 검증으로''의 2단계 구조는 계산과학 분야의
\chg{검증과 확인}(Verification and Validation) 표준에 부합한다.


%% -------------------------------------------------------------------
\chg{이상의 검증 프레임워크를 실제로 수행하기 위한 시뮬레이션의 구체적 구성, 곧 표본 크기와 분포 매개변수, 난수 생성 방법, 변수별 시드 배정 규칙과 그 독립성 검증은 다음 \secref{sec:exp_design}에서 통합적으로 기술한다.}
%% [C3.6/일원화] 구 3.1.2 시드 설계(itemize·그림 fig:seed-bands)는 3.2.3 시드 체계(표 tab:seed_system)·3.2.7 독립성 검증과 중복 → 삭제, 3.2로 일원화(사용자 지시: 분량 무관 중복 삭제).


%% -------------------------------------------------------------------
%% [C3.6/일원화] 구 3.1.3 시뮬레이션 매개변수(itemize·μ=3.0/Lomax 도출)는 3.2.1 tab:sim_parameters(설정근거 포함)·
%% 3.2.5 난수표 각주·3.2.6 분포 설정과 중복 → 삭제, 3.2로 일원화(사용자 지시: 분량 무관 중복 삭제).


%% [C3.3] 구 3.1.4 ``장 요약'' 절 삭제 — 심사 지적(한윤영12·2-B): 절 단위 요약 불필요,
%% 동시에 본 절에 선등장하던 4장 결과수치(C1~C3 평균 효율비·Cohen's d·단축률)도 함께 제거(C4.5 정합).
   % 3.1 연구 설계 (구 2.12)
\section{\vrev{시뮬레이션 설계}}
\label{sec:exp_design}
%% ===================================================================

%% -------------------------------------------------------------------
\subsection{\vrev{몬테카를로 시뮬레이션 구성}}
\label{subsec:mc_config}
%% -------------------------------------------------------------------

본 연구의 시뮬레이션은 조건당 $n = 10{,}000$개의 독립 표본을 생성하는
몬테카를로(Monte Carlo) 방법~\citep{rubinstein2016simulation}에 기반한다.
시뮬레이션의 핵심 매개변수는 다음과 같이
기본 복구시간 $R_0 = 6.0$h, 최대 복구시간 $R_{\max} = 24.0$h($= 4R_0$, MTPD 기반 상한),
시그모이드 곡률 $\kappa = 1.2$이다.
각 표본에 대해 입력 변수를 독립적으로 생성하고,
식~(2.10)에 따라 $\DRTO$ 값과
효율성 비율 $\eta = \DRTO / R_0$를 산출한다.

\tabref{tab:sim_parameters}\refpo{tab:sim_parameters}{은}{는} 시뮬레이션의 전체 매개변수를 요약한 것이며, 각 모델 변수의 전체 정의는 2.7.3(변수 및 매개변수 정의)의 [표 2-11]에 정리되어 있다.

\begin{table}[htbp]
  \centering
  \caption{시뮬레이션 매개변수 요약}
  \label{tab:sim_parameters}
  \small
  \begin{tabularx}{\textwidth}{l c c X}
    \toprule
    \textbf{매개변수} & \textbf{기호} & \textbf{값/범위} & \textbf{설정 근거} \\
    \midrule
    기준선 RTO & $R_0$ & 6.0\,h & \nrev{ISO 22301:2019} BIA 기반 표준 설정값 \\
    DRTO 상한 & $R_{\max}$ & 24.0\,h & MTPD 기반 상한 ($= 4R_0$) \\
    곡률 매개변수 & $\kappa$ & 1.2 & 다기준 복합 점수 최대화 기준 최적값 \\
    표본 크기 & $n$ & 10{,}000 & 조건당, CVaR$_{99}$ 안정 추정 기준 \\
    관측성 원시 점수 & $O_{\text{raw}}$ & $\mathrm{Uniform}(0,1)$ & 연속 균등 분포 \\
    관측성 가중치 & $k_O$ & $\mathrm{Uniform}(0,1)$ & 연속 균등 분포 \\
    불확실성 (C2) & $U_{\text{raw}}^{(G)}$ & $N(3,\;1)$ & 경꼬리, $E[U] = 3.0$ \\
    불확실성 (C3) & $U_{\text{raw}}^{(P)}$ & $6 \cdot \mathrm{Pareto}(\alpha\!=\!3)$ & 파레토(중꼬리), $E[U] = 3.0$ \\
    \bottomrule
  \end{tabularx}
\end{table}


\chg{$n = 10{,}000$의 표본 크기는} 다음의 세 가지 기준에 의해 결정되었다.
첫째, 평균 $\bar{\eta}$의 표준오차가 $\mathrm{SE} = \mathrm{SD}/\sqrt{n}$으로서,
C3의 $\eta$ 표준편차 $\mathrm{SD}(\eta_{\text{C3}}) = 0.791$에서도 $\mathrm{SE} = 0.008$에 불과하여
소수점 둘째 자리까지의 정밀도가 보장된다.
둘째, CVaR$_{99}$의 산출에 상위 $1\%$에 해당하는 100개 이상의 표본이
필요하며, $n = 10{,}000$은 이 요구를 충족한다.
셋째, 부분 표본 분석에서 $n = 5{,}000$ 이상일 때 핵심 통계량이
소수점 셋째 자리까지 수렴함을 사전 확인하였으며 4.6(수렴 진단 및 재현성)에서 그 상세 결과를 알 수 있다.


%% -------------------------------------------------------------------
\subsection{\vrev{난수 생성 방법}}
\label{subsec:rng_method}
%% -------------------------------------------------------------------

본 시뮬레이션의 난수 생성은 [표 3-3]의 시뮬레이션 매개변수 요약 기준에 따라 설계되었으며, 난수 생성기의 상세 규격은 [표 3-4]에 정리하였다.

\begin{table}[htbp]
  \centering
  \caption{난수 생성 방법 규격}
  \label{tab:rng_spec}
  \small
  \begin{tabularx}{\textwidth}{l X}
    \toprule
    \textbf{항목} & \textbf{규격 및 근거} \\
    \midrule
    난수 생성기 & PCG64 (Permuted Congruential Generator, 64-bit). NumPy 1.17+의 \texttt{default\_rng()} 사용.
      기존 Mersenne Twister(MT19937) 대비 통계적 품질·속도·시드 공간 우수. \\
    \m{주기(Period)} & $2^{128}$ (MT19937의 $2^{19937}$보다 작으나,
      $n = 10{,}000$ 수준의 시뮬레이션에서 주기 소진 위험 없음) \\
    시드 지정 & 변수별 개별 시드를 명시적으로 지정하여 \texttt{default\_rng(seed)} 호출.
      암묵적 시스템 시각 기반 시드를 사용하지 않음. (명시 미지정 시 시스템 시각 기반 자동 할당으로 재현 불가) \\
    재현성 보장 & 동일 시드$\cdot$동일 알고리즘$\cdot$동일 호출 순서에서 비트 수준 동일 결과 보장.
      NumPy 버전 간 호환성을 위해 \texttt{Generator} API \m{사용(legacy \texttt{RandomState} 미사용)} \\
    분포 생성 & 균등분포: \texttt{rng.uniform(0, 1)},
      정규분포: \texttt{rng.normal(3, 1)},
      파레토분포: \tnew{\texttt{6 * rng.pareto(3)}} ($E[U] = 3.0$ 보장) \\
    \bottomrule
  \end{tabularx}
\end{table}

\chg{여기서 세 가지 설계 근거는 다음과 같다. 시드 지정과 관련하여, 시드를 명시하지 않으면 많은 난수 생성기가 현재 시스템 시각(system clock)을 시드로 자동 할당하여 실행 시점마다 결과가 달라지므로, 본 연구는 모든 확률 스트림에 고정 시드를 명시적으로 지정하였다. 재현성과 관련하여, NumPy의 구 API인 \texttt{RandomState}는 Mersenne Twister 엔진을 사용하여 버전 갱신 시 동일 시드에서도 난수열이 달라질 수 있는 반면, 신 API인 \texttt{Generator}(NumPy 1.17 이상)는 버전 간 비트 수준 재현성을 보장하므로 이를 채택하였다. 분포 생성과 관련하여, NumPy의 \texttt{Generator.pareto(a)}는 Lomax(파레토 II형) 분포를 반환하여 $E[\mathrm{Lomax}(\alpha=3)] = 1/(\alpha-1) = 0.5$이므로 6을 곱한 \texttt{6 * rng.pareto(3)}의 기댓값이 $3.0$이 된다.}


%% -------------------------------------------------------------------
\subsection{\vrev{비중첩 대역 시드 체계}}
\label{subsec:seed_system}
%% -------------------------------------------------------------------

시뮬레이션의 완전한 재현성을 보장하면서 변수 간 통계적 독립성을
유지하기 위해, 마스터 시드 $s = 42$에서 변수별 시드를 비중첩 대역(non-overlapping band)으로
파생하는 체계를 설계하였다. \tabref{tab:seed_system}\refpo{tab:seed_system}{은}{는}
각 변수의 시드 값과 대역 범위를 정리한 것이다.

\begin{table}[htbp]
  \centering
  \caption{비중첩 대역 시드 \m{체계(마스터 시드 $s = 42$)}}
  \label{tab:seed_system}
  \begin{tabular}{llcc}
    \toprule
    변수 & 기호 & 시드 값 & 대역 범위 \\
    \midrule
    관측성 원시 점수 & $O_{\text{raw}}$ & 42 & $[0\text{--}9999]$ \\
    관측성 가중치 & $k_O$ & 20{,}042 & $[20000\text{--}29999]$ \\
    \m{불확실성(C2, Gaussian)} & $U^{(G)}$ & 30{,}042 & $[30000\text{--}39999]$ \\
    \m{불확실성(C3, Pareto)} & $U^{(P)}$ & 40{,}042 & $[40000\text{--}49999]$ \\
    \bottomrule
  \end{tabular}
\end{table}

각 대역은 $10{,}000$개의 고유 시드를 수용하며, 대역 간 간격이 충분하여
난수 생성기(NumPy PCG64)의 상태 간 교차 오염을 구조적으로
방지한다. 마스터 시드를 변경하면 모든 변수의 시드가 동일한 오프셋 규칙으로
갱신되므로, 다중 시드 검증($s = 0, 1, \ldots, 9{,}999$) 시에도
일관된 독립성이 유지된다.
본 시드 체계의 독립성은 \pdferrR{표~\ref{tab:independence}}에서 실증적으로 검증되며 H6을 충족한다.


%% -------------------------------------------------------------------
\subsection{\padd{재현성 자료 공개}}
\label{subsec:reproducibility-release}
%% -------------------------------------------------------------------

\padd{본 연구의 후속 \chg{검증과 확장} 가능성을 확보하기 위해, 시뮬레이션 코드와
데이터를 GitHub 공개 저장소에 공개한다. \chg{저장소의 공개 주소는 \texttt{https://github.com/WindoorsLEE/drto-dynamic-rto}이며, 코드는 MIT 라이선스로, 데이터는 CC BY 4.0 라이선스로 배포한다.} \chg{공개 자료는 네 종류로 구성된다. 첫째, 시뮬레이션 코드(Python 3.12와 NumPy 2.4 기반)는 C0에서 C3까지의 조건별 몬테카를로 생성과 분할 회귀 및 CVaR 산출, 그리고 $10^8$회 다중 시드 강건성 검증 스크립트를 포함한다. 둘째, 원시 난수열(\texttt{raw\_data/})은 비중첩 대역 시드 체계로 생성된 $O_{\text{raw}}$, $k_O$, $U^{(G)}$, $U^{(P)}$ 4종을 표본 크기 $n = 10{,}000$의 CSV로 담는다. 셋째, 계산 결과(\texttt{calc\_data/})는 조건별 $\DRTO$ 값과 평균, 표준편차, 백분위 등 기준값을 정리한 \texttt{statistics.json}으로 이루어진다. 넷째, LaTeX 소스는 본 논문 본문과 부록을 빌드 가능한 형태로 제공한다.}
재현 절차(\texttt{README.md} 참조)는 저장소 클론 후 \texttt{make all} 실행 시
본 연구의 모든 수치 결과(조건별 \chg{평균, CDF, 회귀계수} 등)가 비트 수준에서
동일하게 재현된다.}


%% -------------------------------------------------------------------
\subsection{\vrev{변수 간 독립성 검증}}
\label{subsec:independence}
%% -------------------------------------------------------------------

비중첩 대역 시드 체계의 유효성을 확인하기 위해, 생성된 $n = 10{,}000$개
표본에서 변수 쌍 간 피어슨 상관계수를 산출하였다.
\tabref{tab:independence}에 제시된 바와 같이, 모든 변수 쌍의
상관계수 절댓값이 $|r| < 0.02$로 나타나 실질적 독립성이 확인되었다. \chg{여기서 기준값 $0.02$는 표본 크기 $n = 10{,}000$에서 귀무가설 $H_0\!: \rho = 0$ 하의 상관계수 표준오차가 $\mathrm{SE}(r) \approx 1/\sqrt{n} = 0.01$이고 95\% 신뢰구간이 $\pm 1.96 \times 0.01 \approx \pm 0.020$이라는 데 근거하며, 따라서 $|r| < 0.02$는 관측된 상관이 유의수준 $\alpha = 0.05$에서 영과 통계적으로 구별되지 않음을 의미한다.}

\begin{table}[htbp]
  \centering
  \caption{입력 변수 간 피어슨 상관계수 ($n = 10{,}000$)}
  \label{tab:independence}
  \begin{tabular}{lcc}
    \toprule
    변수 쌍 & 상관계수 $r$ & 독립성 판정 \\
    \midrule
    $k_O$ -- $O_{\text{raw}}$ & $0.016$ & $|r| < 0.02$ \\
    $O_{\text{raw}}$ -- $U^{(G)}$ & $0.009$ & $|r| < 0.02$ \\
    $O_{\text{raw}}$ -- $U^{(P)}$ & $-0.001$ & $|r| < 0.02$ \\
    \bottomrule
  \end{tabular}
\end{table}


%% -------------------------------------------------------------------
\subsection{\vrev{분포 설정 및 \imod{공정 비교}{공정한 비교} 설계}}
\label{subsec:dist_settings}
%% -------------------------------------------------------------------

C2 조건과 C3 조건의 불확실성 분포는 평균이 동일하되 꼬리 특성이
상이하도록 설계하였다. 이를 통해 분포 형태(shape)가 $\DRTO$에 미치는
고유한 효과를 분리하여 관측할 수 있다.

\chg{C2 조건의 가우시안 불확실성은} $U_{\text{raw}} \sim N(\mu=3,\;\sigma=1)$에서 추출한다.
시그모이드 입력 $z_U = \kappa \cdot k_U \cdot R_0 \cdot U_{\text{raw}} / (R_{\max} - R_0)$로
변환되며, $E[U_{\text{raw}}] = 3.0$의 기대값을 가진다. 시뮬레이션에서
관측된 통계량은 $\bar{U}_{\text{C2}} = 3.003$, $\mathrm{SD} = 0.996$,
$P_{99} = 5.333$이다. 가우시안 분포의 특성상 극단값이 평균에서
크게 이탈하지 않으며, $P_{99}/\mu$ 비율은 $1.78$에 불과하다. \chg{여기서 $P_{99}/\mu$ 비율은 99번째 백분위수와 모평균의 비율로서 분포의 우측 꼬리가 중심으로부터 얼마나 이탈하는지를 나타내며, 본 시뮬레이션에서 $\mu$는 추정값이 아니라 설계 파라미터 $E[U_{\text{raw}}] = 3.0$이므로 정확한 모집단 평균을 사용한다.}

\chg{C3 조건의 파레토 불확실성은} $U_{\text{raw}} \sim 6 \cdot \mathrm{Pareto}(\alpha=3)$에서 추출한다.
형상모수 $\alpha = 3$은 2차 모멘트가 존재하되($\alpha > 2$),
파레토(Heavy-Tail) 특성을 보유하는 분포를 생성한다.
시뮬레이션에서 관측된 통계량은 $\bar{U}_{\text{C3}} = 3.019$,
$\mathrm{SD} = 5.303$, $P_{99} = 23.402$이다.
$P_{99}/\mu$ 비율은 $23.402 / 3.0 = 7.80$으로서,
\textbf{가우시안의 $1.78$ 대비 $4.39$배}에 달하며,
극단 사건의 발생 빈도와 규모가 현저히 크다는 것을 확인할 수 있다.

두 분포의 평균 $\bar{U}_{\text{raw}} \approx 3.0$이 실질적으로
동일하므로, C2와 C3 간 $\eta$ 차이는 불확실성의 크기가 아닌
분포 형태, \chg{특히 우측 꼬리의 중량에} 기인하는 것으로 해석할 수 있다.

\chg{이상에서 제시한 연구 설계와 시뮬레이션 방법에 따라 수행한 실증 분석의 결과는 다음 장에서 조건별 기술통계, 회귀 및 분할 회귀, 강건성 검증의 순으로 제시한다.}
        % 3.2 시뮬레이션 설계 (구 3.1)
